\slajdid{02-s051}
\begin{frame}[label=02-s051]
\gusttekst
\begin{columns}[c,onlytextwidth]
\begin{column}{.63\textwidth}
\centering
\begin{tikzpicture}[circuit logic US,DEoznaka,x=1cm,y=1cm,
 gate/.style={draw,logic gate inputs=nn,minimum width=12mm,minimum height=9mm,scale=.6}]
\node[and gate US,gate] (u) at (2,1.1) {};
\node[and gate US,gate] (d) at (2,-.1) {};
\node[or gate US,gate] (f) at (3.8,.5) {};
\node[not gate US,draw,minimum width=8mm,minimum height=8mm,scale=.6] (i) at (.6,.55) {};
\draw[DEvod] (u.input 1)--(-.5,0 |- u.input 1) node[left] {$C$};
\draw[DEvod] (i.input)--(-.5,.55) node[left] {$B$};
\draw[DEvod] (i.output)--(1.25,.55)|-(u.input 2);
\draw[DEvod] (-.2,.55)|-(d.input 1);
\draw[DEvod] (d.input 2)--(-.5,0 |- d.input 2) node[left] {$A$};
\draw[DEvod] (u.output)--(2.8,1.1)|-(f.input 1);
\draw[DEvod] (d.output)--(2.8,-.1)|-(f.input 2);
\draw[DEvod] (f.output)--++(.4,0) node[right] {$F$};
\end{tikzpicture}
\par
\raggedright
Da bi bolje uočili pojavu, ovaj put ćemo smatrati da su sva logička kola idealno brza, osim invertora koji ima kašnjenje \(t_p\). Neka je stanje signala na ulazu \(C=B=A=1\). Za te ulazne signale \(F=1\). Posmatrajmo promenu samo signala B koji se u nekom trenutku promeni sa logičke jedinice na logičku nulu. Za stacionarno stanje signala \(C=A=1\ i\ B=0\) funkcija takođe ima vrednost \(F=1\). Znači bez obzira na promenu stanje izlaza treba da je logička jedinica.
\par\centering ALI
\end{column}
\begin{column}{.33\textwidth}\centering
\begin{tikzpicture}[x=1cm,y=1cm,font=\fontsize{8}{9}\selectfont]
\foreach \r/\lab in {0/A,1/B,2/{\bar B},3/C,4/A,5/{BA},6/{C\bar B},7/F}{
 \pgfmathsetmacro{\yy}{-.79*\r}
 \draw[DEstrelica] (0,\yy-.06)--(0,\yy+.60) node[left] {$\lab$};
 \draw[DEstrelica] (-.1,\yy)--(3.8,\yy) node[below] {$t$};
}
\foreach \r in {0,3,4}{
 \pgfmathsetmacro{\yy}{-.79*\r}
 \draw[DEvod] (0,\yy+.44)--(3.35,\yy+.44);
}
\foreach \r in {1,5}{
 \pgfmathsetmacro{\yy}{-.79*\r}
 \draw[DEvod] (0,\yy+.44)--(.95,\yy+.44)--(1.06,\yy+.08)--(3.35,\yy+.08);
}
\foreach \r in {2,6}{
 \pgfmathsetmacro{\yy}{-.79*\r}
 \draw[DEvod] (0,\yy+.08)--(1.76,\yy+.08)--(1.88,\yy+.44)--(3.35,\yy+.44);
}
\draw[DEvod] (0,-5.09)--(.95,-5.09)--(1.06,-5.45)--(1.76,-5.45)--(1.88,-5.09)--(3.35,-5.09);
\draw[DEvod,dashed] (1,.58)--(1,-5.62);
\draw[DEvod,dashed] (1.82,.58)--(1.82,-5.62);
\draw[DEvod,{Latex[length=1.1mm]}-{Latex[length=1.1mm]}] (1,-1.70)--(1.82,-1.70)
 node[midway,below,inner sep=1pt] {$t_p$};
\node[font=\fontsize{8}{9}\selectfont,inner sep=1pt] at (2.7,-5.95) {lažna nula};
\end{tikzpicture}

\end{column}
\end{columns}
\end{frame}
